\documentclass{amsart}
% For dealing with references we use the comment environment
\usepackage{verbatim}
\newenvironment{reference}{\comment}{\endcomment}
%\newenvironment{reference}{}{}
\newenvironment{slogan}{\comment}{\endcomment}
\newenvironment{history}{\comment}{\endcomment}

% For commutative diagrams we use Xy-pic
\usepackage[all]{xy}

% We use 2cell for 2-commutative diagrams.
\xyoption{2cell}
\UseAllTwocells

% We use multicol for the list of chapters between chapters
\usepackage{multicol}

% This is generally recommended for better output
\usepackage{lmodern}
\usepackage[T1]{fontenc}

% For cross-file-references

% Package for hypertext links:
\usepackage{hyperref}

% For any local file, say "hello.tex" you want to link to please

% Theorem environments.
%
\theoremstyle{plain}
\newtheorem{theorem}[subsection]{Theorem}
\newtheorem{proposition}[subsection]{Proposition}
\newtheorem{lemma}[subsection]{Lemma}

\theoremstyle{definition}
\newtheorem{definition}[subsection]{Definition}
\newtheorem{example}[subsection]{Example}
\newtheorem{exercise}[subsection]{Exercise}
\newtheorem{situation}[subsection]{Situation}

\theoremstyle{remark}
\newtheorem{remark}[subsection]{Remark}
\newtheorem{remarks}[subsection]{Remarks}

\numberwithin{equation}{subsection}

% Macros
%
\def\lim{\mathop{\mathrm{lim}}\nolimits}
\def\colim{\mathop{\mathrm{colim}}\nolimits}
\def\Spec{\mathop{\mathrm{Spec}}}
\def\Hom{\mathop{\mathrm{Hom}}\nolimits}
\def\Ext{\mathop{\mathrm{Ext}}\nolimits}
\def\SheafHom{\mathop{\mathcal{H}\!\mathit{om}}\nolimits}
\def\SheafExt{\mathop{\mathcal{E}\!\mathit{xt}}\nolimits}
\def\Sch{\mathit{Sch}}
\def\Mor{\mathop{\mathrm{Mor}}\nolimits}
\def\Ob{\mathop{\mathrm{Ob}}\nolimits}
\def\Sh{\mathop{\mathit{Sh}}\nolimits}
\def\NL{\mathop{N\!L}\nolimits}
\def\CH{\mathop{\mathrm{CH}}\nolimits}
\def\proetale{{pro\text{-}\acute{e}tale}}
\def\etale{{\acute{e}tale}}
\def\QCoh{\mathit{QCoh}}
\def\Ker{\mathop{\mathrm{Ker}}}
\def\Im{\mathop{\mathrm{Im}}}
\def\Coker{\mathop{\mathrm{Coker}}}
\def\Coim{\mathop{\mathrm{Coim}}}

% Boxtimes
%
\DeclareMathSymbol{\boxtimes}{\mathbin}{AMSa}{"02}

%
% Macros for moduli stacks/spaces
%
\def\QCohstack{\mathcal{QC}\!\mathit{oh}}
\def\Cohstack{\mathcal{C}\!\mathit{oh}}
\def\Spacesstack{\mathcal{S}\!\mathit{paces}}
\def\Quotfunctor{\mathrm{Quot}}
\def\Hilbfunctor{\mathrm{Hilb}}
\def\Curvesstack{\mathcal{C}\!\mathit{urves}}
\def\Polarizedstack{\mathcal{P}\!\mathit{olarized}}
\def\Complexesstack{\mathcal{C}\!\mathit{omplexes}}
% \Pic is the operator that assigns to X its picard group, usage \Pic(X)
% \Picardstack_{X/B} denotes the Picard stack of X over B
% \Picardfunctor_{X/B} denotes the Picard functor of X over B
\def\Pic{\mathop{\mathrm{Pic}}\nolimits}
\def\Picardstack{\mathcal{P}\!\mathit{ic}}
\def\Picardfunctor{\mathrm{Pic}}
\def\Deformationcategory{\mathcal{D}\!\mathit{ef}}

\setlength{\emergencystretch}{3em}
\begin{document}
\title[Stacks Project, Tag 019G]{706 --- Dold-Kan equivalence for an arbitrary abelian category}
\maketitle
\noindent Copyright (C) 2005--2025 Johan de Jong. Stacks Project authors. Source: \href{https://stacks.math.columbia.edu/tag/019G}{Tag 019G}.

\noindent Editorial difficulty note (not author text). Difficulty H2: The author constructs an explicit inverse indexed by surjective simplex maps, verifies functoriality by a complete case analysis, and proves both natural inverse isomorphisms using normalized splitting. This full constructive equivalence is substantially beyond ordinary qualification homological algebra..

\noindent Editorial provenance note (not author text). History: excerpt from revision \texttt{a0444\allowbreak 6e57e\allowbreak c1fbc\allowbreak 252a8\allowbreak 71afc\allowbreak ec775\allowbreak 2fb28\allowbreak 07b14}, simplicial.tex, lines 4044--4273. Reference presentation was adapted; original-author context and core support blocks were retained or added. The mathematical wording is unchanged. Licensed under GNU FDL 1.2 or later; no invariant sections or cover texts. See ../licenses/COPYING. Context blocks reproduce author definitions and supporting results; they are not additional counted problems.

\begin{definition}
\label{definition-split}
Let $\mathcal{C}$ be a category which admits finite nonempty coproducts.
We say a simplicial object $U$ of $\mathcal{C}$ is {\it split}
if there exist subobjects $N(U_m)$ of $U_m$, $m \geq 0$
with the property that
\begin{equation}
\label{equation-splitting}
\coprod\nolimits_{\varphi : [n] \to [m]\text{ surjective}}
N(U_m)
\longrightarrow
U_n
\end{equation}
is an isomorphism for all $n \geq 0$. If $U$ is an $r$-truncated
simplicial object of $\mathcal{C}$ then we say $U$ is {\it split}
if there exist subobjects $N(U_m)$ of $U_m$, $r \geq m \geq 0$
with the property that (\ref{equation-splitting})
is an isomorphism for $r \geq n \geq 0$.
\end{definition}

\noindent
There is a second chain complex we can associate to
a simplicial object of $\mathcal{A}$. Recall that by
Lemma \ref{lemma-splitting-abelian-category}
any simplicial object $U$ of $\mathcal{A}$
is canonically split with
$N(U_m) = \bigcap_{i = 0}^{m - 1} \Ker(d^m_i)$.
We define the {\it normalized chain complex $N(U)$}
to be the chain complex
$$
\ldots \to N(U_2) \to N(U_1) \to N(U_0) \to 0 \to 0 \to \ldots
$$
with boundary map $d_n : N(U_n) \to N(U_{n - 1})$ given
by the restriction of $(-1)^nd^n_n$ to the direct summand
$N(U_n)$ of $U_n$. Note that Lemma \href{https://stacks.math.columbia.edu/tag/017X}{lemma N d in N (Tag 017X)}
implies that $d^n_n(N(U_n)) \subset N(U_{n - 1})$.
It is a complex because
$d^n_n \circ d^{n + 1}_{n + 1} = d^n_n \circ d^{n + 1}_n$
and $d^{n + 1}_n$ is zero on $N(U_{n + 1})$ by definition.
Thus we obtain a second functor
$$
N : \text{Simp}(\mathcal{A}) \longrightarrow \text{Ch}_{\geq 0}(\mathcal{A}).
$$


\begin{lemma}
\label{lemma-splitting-simplicial-groups}
Let $U$ be a simplicial abelian group.
Then $U$ has a splitting obtained by taking $N(U_0) = U_0$ and
for $m \geq 1$ taking
$$
N(U_m) = \bigcap\nolimits_{i = 0}^{m - 1} \Ker(d^m_i).
$$
Moreover, this splitting is functorial on the category
of simplicial abelian groups.
\end{lemma}

\begin{proof}
By induction on $n$ we will show that the choice of $N(U_m)$
in the lemma guarantees that (\ref{equation-splitting}) is
an isomorphism for $m \leq n$. This is clear for $n = 0$.
In the rest of this proof we are going to
drop the superscripts from the maps $d_i$ and $s_i$ in order
to improve readability. We will also repeatedly use the relations
from Remark \href{https://stacks.math.columbia.edu/tag/016C}{remark relations (Tag 016C)}.

\medskip\noindent
First we make a general remark.
For $0 \leq i \leq m$ and $z \in U_m$ we have
$d_i(s_i(z)) = z$. Hence we can write
any $x \in U_{m + 1}$ uniquely as
$x = x' + x''$ with $d_i(x') = 0$
and $x'' \in \Im(s_i)$
by taking $x' = (x - s_i(d_i(x)))$ and
$x'' = s_i(d_i(x))$. Moreover, the element
$z \in U_m$ such that $x'' = s_i(z)$
is unique because $s_i$ is injective.

\medskip\noindent
Here is a procedure for decomposing any $x \in U_{n + 1}$.
First, write $x = x_0 + s_0(z_0)$ with $d_0(x_0) = 0$.
Next, write $x_0 = x_1 + s_1(z_1)$ with
$d_1(x_1) = 0$. Continue like this to get
\begin{eqnarray*}
x & = & x_0 + s_0(z_0), \\
x_0 & = & x_1 + s_1(z_1), \\
x_1 & = & x_2 + s_2(z_2), \\
\ldots & \ldots & \ldots \\
x_{n - 1} & = & x_n + s_n(z_n)
\end{eqnarray*}
where $d_i(x_i) = 0$ for all $i = n, \ldots, 0$.
By our general remark above all of the $x_i$
and $z_i$ are determined uniquely by $x$.
We claim that
$x_i \in
\Ker(d_0) \cap
\Ker(d_1) \cap
\ldots \cap
\Ker(d_i)$
and
$z_i \in
\Ker(d_0) \cap
\ldots \cap
\Ker(d_{i - 1})$
for $i = n, \ldots, 0$.
Here and in the following
an empty intersection of kernels indicates
the whole space; i.e.,
the notation
$z_0 \in \Ker(d_0) \cap
\ldots \cap
\Ker(d_{i - 1})$
when $i = 0$ means $z_0 \in U_n$ with no restriction.

\medskip\noindent
We prove this by ascending induction on $i$.
It is clear for $i = 0$ by construction of $x_0$ and $z_0$.
Let us prove it for $0 < i \leq n$ assuming the result for $i - 1$.
First of all we have $d_i(x_i) = 0$ by construction.
So pick a $j$ with $0 \leq j < i$. We have
$d_j(x_{i - 1}) = 0$ by induction. Hence
$$
0 = d_j(x_{i - 1})
= d_j(x_i) + d_j(s_i(z_i))
= d_j(x_i) + s_{i - 1}(d_j(z_i)).
$$
The last equality by the relations of Remark \href{https://stacks.math.columbia.edu/tag/016C}{remark relations (Tag 016C)}.
These relations also imply that
$d_{i - 1}(d_j(x_i)) = d_j(d_i(x_i)) = 0$
because $d_i(x_i)= 0$ by construction.
Then the uniqueness in the general remark above shows the equality
$0 = x' + x'' = d_j(x_i) + s_{i - 1}(d_j(z_i))$
can only hold if both terms are zero. We conclude that
$d_j(x_i) = 0$ and by injectivity of $s_{i - 1}$ we also
conclude that $d_j(z_i) = 0$. This proves the claim.

\medskip\noindent
The claim implies we can uniquely write
$$
x = s_0(z_0) + s_1(z_1) + \ldots + s_n(z_n) + x_n
$$
with $x_n \in N(U_{n + 1})$ and
$z_i \in \Ker(d_0) \cap \ldots \cap \Ker(d_{i - 1})$.
We can reformulate this as saying that we have found a direct
sum decomposition
$$
U_{n + 1}
=
N(U_{n + 1})
\oplus
\bigoplus\nolimits_{i = 0}^{i = n}
s_i\Big(\Ker(d_0) \cap \ldots \cap \Ker(d_{i - 1})\Big)
$$
with the property that
$$
\Ker(d_0) \cap \ldots \cap \Ker(d_j)
=
N(U_{n + 1}) \oplus
\bigoplus\nolimits_{i = j + 1}^{i = n}
s_i\Big(\Ker(d_0) \cap \ldots \cap \Ker(d_{i - 1})\Big)
$$
for $j = 0, \ldots, n$.
The result follows from this statement as follows.
Each of the $z_i$ in the expression for $x$
can be written uniquely as
$$
z_i = s_i(z'_{i, i}) + \ldots + s_{n - 1}(z'_{i, n - 1}) + z_{i, 0}
$$
with $z_{i, 0} \in N(U_n)$ and
$z'_{i, j} \in \Ker(d_0) \cap \ldots \cap \Ker(d_{j - 1})$.
The first few steps in the decomposition of $z_i$ are zero because
$z_i$ already is in the kernel of $d_0, \ldots, d_i$.
This in turn uniquely gives
$$
x = x_n + s_0(z_{0, 0}) + s_1(z_{1, 0}) + \ldots + s_n(z_{n, 0}) +
\sum\nolimits_{0 \leq i \leq j \leq n - 1} s_i(s_j(z'_{i, j})).
$$
Continuing in this fashion we see that we in the end obtain
a decomposition of $x$ as a sum of terms
of the form
$$
s_{i_1} s_{i_2} \ldots s_{i_k} (z)
$$
with $0 \leq i_1 \leq i_2 \leq \ldots \leq i_k \leq n - k + 1$ and
$z \in N(U_{n + 1 - k})$. This is exactly the required
decomposition, because any surjective map $[n + 1] \to [n + 1 - k]$
can be uniquely expressed in the form
$$
\sigma^{n - k}_{i_k} \ldots \sigma^{n - 1}_{i_2} \sigma^n_{i_1}
$$
with $0 \leq i_1 \leq i_2 \leq \ldots \leq i_k \leq n - k + 1$.
\end{proof}


\begin{lemma}
\label{lemma-splitting-abelian-category}
Let $\mathcal{A}$ be an abelian category.
Let $U$ be a simplicial object in $\mathcal{A}$.
Then $U$ has a splitting obtained by taking $N(U_0) = U_0$ and
for $m \geq 1$ taking
$$
N(U_m) = \bigcap\nolimits_{i = 0}^{m - 1} \Ker(d^m_i).
$$
Moreover, this splitting is functorial on the category of
simplicial objects of $\mathcal{A}$.
\end{lemma}

\begin{proof}
For any object $A$ of $\mathcal{A}$ we obtain
a simplicial abelian group $\Mor_\mathcal{A}(A, U)$.
Each of these are canonically split by Lemma
\ref{lemma-splitting-simplicial-groups}. Moreover,
$$
N(\Mor_\mathcal{A}(A, U_m)) =
\bigcap\nolimits_{i = 0}^{m - 1} \Ker(d^m_i) =
\Mor_\mathcal{A}(A, N(U_m)).
$$
Hence we see that the morphism (\ref{equation-splitting})
becomes an isomorphism after applying the functor
$\Mor_\mathcal{A}(A, -)$ for any object of $\mathcal{A}$.
Hence it is an isomorphism by the Yoneda lemma.
\end{proof}


\begin{lemma}
\label{lemma-decompose-associated-complexes}
Let $\mathcal{A}$ be an abelian category.
Let $U$ be a simplicial object of $\mathcal{A}$.
The canonical morphism of chain complexes
$N(U) \to s(U)$ is split. In fact,
$$
s(U) = N(U) \oplus D(U)
$$
for some complex $D(U)$. The construction $U \mapsto D(U)$
is functorial.
\end{lemma}

\begin{proof}
Define $D(U)_n$ to be the image of
$$
\bigoplus\nolimits_{\varphi : [n] \to [m]\text{ surjective}, \ m < n} N(U_m)
\xrightarrow{\bigoplus U(\varphi)}
U_n
$$
which is a subobject of $U_n$ complementary to
$N(U_n)$ according to Lemma \ref{lemma-splitting-abelian-category} and
Definition \ref{definition-split}. We show that
$D(U)$ is a subcomplex. Pick a surjective
map $\varphi : [n] \to [m]$ with $m < n$ and consider
the composition
$$
N(U_m) \xrightarrow{U(\varphi)} U_n \xrightarrow{d_n} U_{n - 1}.
$$
This composition is the sum of the maps
$$
N(U_m) \xrightarrow{U(\varphi \circ \delta^n_i)} U_{n - 1}
$$
with sign $(-1)^i$, $i = 0, \ldots, n$.

\medskip\noindent
First we will prove by ascending induction on $m$,
$0 \leq m < n - 1$ that all the maps $U(\varphi \circ \delta^n_i)$
map $N(U_m)$ into $D(U)_{n - 1}$. (The case $m = n - 1$ is treated below.)
Whenever the map $\varphi \circ \delta^n_i : [n - 1] \to [m]$
is surjective then the image of $N(U_m)$ under $U(\varphi \circ \delta^n_i)$
is contained in $D(U)_{n - 1}$ by definition.
If $\varphi \circ \delta^n_i : [n - 1] \to [m]$ is not surjective,
set $j = \varphi(i)$ and observe that $i$ is the unique
index whose image under $\varphi$ is $j$. We may write
$\varphi \circ \delta^n_i = \delta^m_j \circ \psi \circ \delta^n_i$
for some $\psi : [n - 1] \to [m - 1]$. Hence
$U(\varphi \circ \delta^n_i) = U(\psi \circ \delta^n_i) \circ d^m_j$ which
is zero on $N(U_m)$ unless $j = m$. If $j = m$, then
$d^m_m(N(U_m)) \subset N(U_{m - 1})$ and hence
$U(\varphi \circ \delta^n_i)(N(U_m)) \subset
U(\psi \circ \delta^n_i)(N(U_{m - 1}))$ and we win
by induction hypothesis.

\medskip\noindent
To finish proving that $D(U)$ is a subcomplex
we still have to deal with the composition
$$
N(U_m) \xrightarrow{U(\varphi)} U_n \xrightarrow{d_n} U_{n - 1}.
$$
in case $m =  n - 1$. In this case $\varphi = \sigma^{n - 1}_j$
for some $0 \leq j \leq n - 1$ and $U(\varphi) = s^{n - 1}_j$.
Thus the composition is given by the sum
$$
\sum (-1)^i d^n_i \circ s^{n - 1}_j
$$
Recall from Remark \href{https://stacks.math.columbia.edu/tag/016C}{remark relations (Tag 016C)} that
$d^n_j \circ s^{n - 1}_j = d^n_{j + 1} \circ s^{n - 1}_j = \text{id}$
and these drop out because the corresponding terms have opposite signs.
The map $d^n_n \circ s^{n - 1}_j$, if $j < n - 1$, is equal to
$s^{n - 2}_j \circ d^{n - 1}_{n - 1}$. Since
$d^{n - 1}_{n - 1}$ maps $N(U_{n - 1})$ into $N(U_{n - 2})$,
we see that the image $d^n_n ( s^{n - 1}_j (N(U_{n - 1}))$
is contained in $s^{n - 2}_j(N(U_{n - 2}))$ which
is contained in $D(U_{n - 1})$ by definition. For all
other combinations of $(i, j)$ we have
either $d^n_i \circ s^{n - 1}_j = s^{n - 2}_{j - 1} \circ d^{n - 1}_i$
(if $i < j$), or
$d^n_i \circ s^{n - 1}_j = s^{n - 2}_j \circ d^{n - 1}_{i - 1}$
(if $n > i > j + 1$) and in these cases the map is zero because
of the definition of $N(U_{n - 1})$.
\end{proof}


\begin{theorem}
\label{theorem-dold-kan}
Let $\mathcal{A}$ be an abelian category.
The functor $N$ induces an equivalence of
categories
$$
N :
\text{Simp}(\mathcal{A})
\longrightarrow
\text{Ch}_{\geq 0}(\mathcal{A})
$$
\end{theorem}

\begin{proof}
We will describe a functor in the reverse direction
inspired by the construction of Lemma \href{https://stacks.math.columbia.edu/tag/0192}{lemma extension (Tag 0192)}
(except that we throw in a sign to get the boundaries
right). Let $A_\bullet$ be a chain complex with boundary maps
$d_{A, n} : A_n \to A_{n - 1}$. For each $n \geq 0$ denote
$$
I_n =
\Big\{
\alpha : [n] \to \{0, 1, 2, \ldots\}
\mid
\Im(\alpha) = [k]\text{ for some }k
\Big\}.
$$
For $\alpha \in I_n$ we denote $k(\alpha)$ the unique
integer such that $\Im(\alpha) = [k]$.
We define a simplicial object $S(A_\bullet)$ as follows:
\begin{enumerate}
\item $S(A_\bullet)_n = \bigoplus_{\alpha \in I_n} A_{k(\alpha)}$, which
we will write as
$\bigoplus_{\alpha \in I_n} A_{k(\alpha)} \cdot \alpha$
to suggest thinking of ``$\alpha$'' as a basis vector for the
summand corresponding to it,
\item given $\varphi : [m] \to [n]$ we define
$S(A_\bullet)(\varphi)$ by its restriction to
the direct summand $A_{k(\alpha)} \cdot \alpha$
of $S(A_\bullet)_n$ as follows
\begin{enumerate}
\item $\alpha \circ \varphi \not \in I_m$ then we set it equal to zero,
\item $\alpha \circ \varphi \in I_m$ but $k(\alpha \circ \varphi)$
not equal to either $k(\alpha)$ or $k(\alpha) - 1$ then we set it
equal to zero as well,
\item if $\alpha \circ \varphi \in I_m$
and $k(\alpha \circ \varphi) = k(\alpha)$ then we use
the identity map to the summand
$A_{k(\alpha \circ \varphi)} \cdot (\alpha \circ \varphi)$
of $S(A_\bullet)_m$, and
\item if $\alpha \circ \varphi \in I_m$
and $k(\alpha \circ \varphi) = k(\alpha) - 1$
then we use $(-1)^{k(\alpha)} d_{A, k(\alpha)}$ to the summand
$A_{k(\alpha \circ \varphi)}\cdot (\alpha \circ \varphi)$
of $S(A_\bullet)_m$.
\end{enumerate}
\end{enumerate}
Let us show that $S(A_\bullet)$ is a simplicial object of $\mathcal{A}$.
To do this, assume we have maps $\varphi : [m] \to [n]$ and
$\psi : [n] \to [p]$. We will show that
$S(A_\bullet)(\varphi) \circ S(A_\bullet)(\psi) =
S(A_\bullet)(\psi \circ \varphi)$. Choose $\beta \in I_p$ and set
$\alpha = \beta \circ \psi$ and $\gamma = \alpha \circ \varphi$
viewed as maps $\alpha : [n] \to \{0, 1, 2, \ldots\}$
and $\gamma : [m] \to \{0, 1, 2, \ldots\}$. Picture
$$
\xymatrix{
[m] \ar[r]_\varphi \ar[d]_\gamma &
[n] \ar[r]_\psi \ar[d]_\alpha &
[p] \ar[d]^\beta \\
\Im(\gamma) \ar[r] &
\Im(\alpha) \ar[r] &
[k(\beta)]
}
$$
We will show that the restriction of the maps
$S(A_\bullet)(\varphi) \circ S(A_\bullet)(\psi)$ and
$S(A_\bullet)(\psi \circ \varphi)$.
to the summand $A_{k(\beta)} \cdot \beta$ agree.
There are several cases to consider
\begin{enumerate}
\item Say $\alpha \not \in I_n$ so the restriction of
$S(A_\bullet)(\psi)$ to $A_{k(\beta)} \cdot \beta$ is zero.
Then either $\gamma \not \in I_m$ or we have
$[k(\gamma)] = \Im(\gamma) \subset \Im(\alpha) \subset [k(\beta)]$
and the subset $\Im(\alpha)$ of $[k(\beta)]$ has a gap so
$k(\gamma) < k(\beta) - 1$. In both cases we see that
the restriction of
$S(A_\bullet)(\psi \circ \varphi)$ to $A_{k(\beta)} \cdot \beta$
is zero as well.
\item Say $\alpha \in I_n$ and $k(\alpha) < k(\beta) - 1$ so the restriction of
$S(A_\bullet)(\psi)$ to $A_{k(\beta)} \cdot \beta$ is zero.
Then either $\gamma \not \in I_m$ or we have
$[k(\gamma)] \subset [k(\alpha)] \subset [k(\beta)]$
and it follows that $k(\gamma) < k(\beta) - 1$. In both cases we see that
the restriction of
$S(A_\bullet)(\psi \circ \varphi)$ to $A_{k(\beta)} \cdot \beta$
is zero as well.
\item Say $\alpha \in I_n$ and $k(\alpha) = k(\beta)$ so the restriction of
$S(A_\bullet)(\psi)$ to $A_{k(\beta)} \cdot \beta$ is
the identity map from $A_{k(\beta)} \cdot \beta$ to
$A_{k(\alpha)} \cdot \alpha$. In this case because $\Im(\alpha) =[k(\beta)]$
the rule describing the restriction of $S(A_\bullet)(\psi \circ \varphi)$
to the summand $A_{k(\beta)} \cdot \beta$ is exactly the same
as the rule describing the restriction of $S(A_\bullet)(\varphi)$
to the summand $A_{k(\alpha)} \cdot \alpha$ and hence agreement holds.
\item Say $\alpha \in I_n$ and $k(\alpha) = k(\beta) - 1$
so the restriction of $S(A_\bullet)(\psi)$ to $A_{k(\beta)} \cdot \beta$
is given by $(-1)^{k(\beta)}d_{A, k(\beta)}$ to $A_{k(\alpha)} \cdot \alpha$.
Subcases
\begin{enumerate}
\item If $\gamma \not \in I_m$, then both the restriction of
$S(A_\bullet)(\psi \circ \varphi)$ to the summand $A_{k(\beta)} \cdot \beta$
and the restriction of $S(A_\bullet)(\varphi)$ to the summand
$A_{k(\alpha)} \cdot \alpha$ are zero and we get agreement.
\item If $\gamma \in I_m$ but $k(\gamma) < k(\alpha) - 1$, then
again both restrictions are zero and we get agreement.
\item If $\gamma \in I_m$ and $k(\gamma) = k(\alpha)$ then
$\Im(\gamma) = \Im(\alpha)$. In this case
the restriction of $S(A_\bullet)(\psi \circ \varphi)$
to the summand $A_{k(\beta)} \cdot \beta$ is given by
$(-1)^{k(\beta)}d_{A, k(\beta)}$ to $A_{k(\gamma)} \cdot \gamma$
and the restriction of $S(A_\bullet)(\varphi)$ to the summand
$A_{k(\alpha)} \cdot \alpha$ is the identity map
$A_{k(\alpha)} \cdot \alpha \to A_{k(\gamma)} \cdot \gamma$.
Hence agreement holds.
\item Finally, if $\gamma \in I_m$ and $k(\gamma) = k(\alpha) - 1$ then
the restriction of $S(A_\bullet)(\varphi)$ to the summand
$A_{k(\alpha)} \cdot \alpha$ is given by $(-1)^{k(\alpha)} d_{A, k(\alpha)}$
as a map
$A_{k(\alpha)} \cdot \alpha \to A_{k(\beta)} \cdot \beta$.
Since $A_\bullet$ is a complex we see that the composition
$A_{k(\beta)} \cdot \beta \to
A_{k(\alpha)} \cdot \alpha \to A_{k(\gamma)} \cdot \gamma$ is zero
which matches what we get for the restriction of
$S(A_\bullet)(\psi \circ \varphi)$
to the summand $A_{k(\beta)} \cdot \beta$
because $k(\gamma) = k(\beta) - 2 < k(\beta) - 1$.
\end{enumerate}
\end{enumerate}
Thus $S(A_\bullet)$ is a simplicial object of $\mathcal{A}$.

\medskip\noindent
Let us construct an isomorphism $A_\bullet \to N(S(A_\bullet))$
functorial in $A_\bullet$. Recall that
$$
S(A_\bullet) =  N(S(A_\bullet)) \oplus D(S(A_\bullet))
$$
as chain complexes by Lemma \ref{lemma-decompose-associated-complexes}.
On the other hand it follows from Remark \href{https://stacks.math.columbia.edu/tag/0FKI}{remark degenerate subcomplex (Tag 0FKI)}
and the construction of $S(A_\bullet)$ that
$$
D(S(A_\bullet))_n =
\bigoplus\nolimits_{\alpha \in I_n,\ k(\alpha) < n} A_{k(\alpha)} \cdot \alpha
\subset
\bigoplus\nolimits_{\alpha \in I_n} A_{k(\alpha)} \cdot \alpha
$$
However, if $\alpha \in I_n$ then we have
$k(\alpha) \geq n \Leftrightarrow \alpha = \text{id}_{[n]} : [n] \to [n]$.
Thus the summand $A_n \cdot \text{id}_{[n]}$ of $S(A_\bullet)_n$
is a complement to the summand $D(S(A_\bullet))_n$.
All the maps $d^n_i : S(A_\bullet)_n \to S(A_\bullet)_n$
restrict to zero on the summand $A_n \cdot \text{id}_{[n]}$
except for $d^n_n$ which produces $(-1)^n d_{A, n}$
from $A_n \cdot \text{id}_{[n]}$ to $A_{n - 1} \cdot \text{id}_{[n - 1]}$.
We conclude that $A_n \cdot \text{id}_{[n]}$ must be
equal to the summand $N(S(A_\bullet))_n$ and moreover
the restriction of the differential
$d_n = \sum (-1)^id^n_i : S(A_\bullet)_n \to S(A_\bullet)_{n - 1}$
to the summand $A_n \cdot \text{id}_{[n]}$ gives what we want!

\medskip\noindent
Finally, we have to show that $S \circ N$ is isomorphic to the
identity functor. Let $U$ be a simplicial object of $\mathcal{A}$.
Then we can define an obvious map
$$
S(N(U))_n = \bigoplus\nolimits_{\alpha \in I_n} N(U)_{k(\alpha)} \cdot \alpha
\longrightarrow
U_n
$$
by using $U(\alpha) : N(U)_{k(\alpha)} \to U_n$ on the summand corresponding
to $\alpha$. By Definition \ref{definition-split} this is an isomorphism.
To finish the proof we have to show that this is compatible with
the maps in the simplicial objects. Thus let $\varphi : [m] \to [n]$
and let $\alpha \in I_n$. Set $\beta = \alpha \circ \varphi$.
Picture
$$
\xymatrix{
[m] \ar[r]_\varphi \ar[d]_\beta &
[n] \ar[d]_\alpha \\
\Im(\beta) \ar[r] &
[k(\alpha)]
}
$$
There are several cases to consider
\begin{enumerate}
\item Say $\beta \not \in I_m$. Then there exists an index
$0 \leq j < k(\alpha)$ with $j \not \in \Im(\alpha \circ \varphi)$
and hence we can choose a factorization
$\alpha \circ \varphi = \delta^{k(\alpha)}_j \circ \psi$
for some $\psi : [m] \to [k(\alpha) - 1]$.
It follows that $U(\varphi)$ is zero on the image of the summand
$N(U)_{k(\alpha)} \cdot \alpha$ because $U(\varphi) \circ U(\alpha) =
U(\alpha \circ \varphi) = U(\psi) \circ d^{k(\alpha)}_j$
is zero on $N(U)_{k(\alpha)}$ by construction of $N$.
This matches our rule for $S(N(U))$ given above.
\item Say $\beta \in I_m$ and $k(\beta) < k(\alpha) - 1$. Here we
argue exactly as in case (1) with $j = k(\alpha) - 1$.
\item Say $\beta \in I_m$ and $k(\beta) = k(\alpha)$. Here
the summand $N(U)_{k(\alpha)} \cdot \alpha$ is mapped by the
identity to the summand $N(U)_{k(\beta)} \cdot \beta$.
This is the same as the effect of $U(\varphi)$ since in
this case $U(\varphi) \circ U(\alpha) = U(\beta)$.
\item Say $\beta \in I_m$ and $k(\beta) = k(\alpha) - 1$.
Here we use the differential $(-1)^{k(\alpha)} d_{N(U), k(\alpha)}$
to map the summand $N(U)_{k(\alpha)} \cdot \alpha$
to the summand $N(U)_{k(\beta)} \cdot \beta$.
On the other hand, since $\Im(\beta) = [k(\beta)]$ in this case
we get $\alpha \circ \varphi = \delta^{k(\alpha)}_{k(\alpha)} \circ \beta$.
Thus we see that $U(\varphi)$ composed with the restriction of
$U(\alpha)$ to $N(U)_{k(\alpha)}$ is equal to
$U(\beta)$ precomposed with $d^{k(\alpha)}_{k(\alpha)}$ restricted to
$N(U)_{k(\alpha)}$.
Since $d_{N(U), k(\alpha)} = \sum (-1)^i d^{k(\alpha)}_i$
and since $d^{k(\alpha)}_i$ restricts to zero on
$N(U)_{k(\alpha)}$ for $i < k(\alpha)$
we see that equality holds.
\end{enumerate}
This finishes the proof of the theorem.
\end{proof}
\end{document}
